\chapter*{Abstrak}
\addcontentsline{toc}{chapter}{Abstrak}

% --- Isi Teks Abstrak Bahasa Indonesia ---
%
%     Tulis isi teks abstrak Anda dengan maksimal 250 kata
%
\newcommand{\abstrak}{%
	Abstrak ditulis sebagai teks yang menyajikan esensi dokumen secara mandiri tanpa merujuk pada tabel, gambar, ataupun daftar pustaka. Komposisi teks diawali dengan latar belakang atau rasional penelitian yang melandasi pentingnya kajian atau pengembangan topik ini diangkat. Setelah urgensi masalah dipaparkan secara singkat, bagian selanjutnya wajib merumuskan tujuan utama yang ingin dicapai dalam penelitian secara jelas. Komponen berikutnya yang sangat krusial adalah penjelasan mengenai metode atau pendekatan pemecahan masalah yang digunakan oleh peneliti. Penjelasan metodologi ini harus mencakup desain penelitian, lokasi, subjek atau sumber data, teknik pengumpulan data, instrumen yang digunakan, serta teknik analisis data yang diterapkan untuk mengolah informasi. Setelah metode diuraikan, struktur beralih pada pemaparan hasil temuan utama yang diperoleh berdasarkan analisis objek kajian secara objektif. Bagian akhir dari tubuh abstrak ditutup dengan simpulan yang menjawab langsung tujuan awal, serta saran operasional yang sesuai bagi pengembangan ilmu maupun kebijakan praktis ke depan. Seluruh uraian tersebut wajib dirangkai secara padat, lugas, dan sistematis dengan menggunakan spasi tunggal dan panjang total tidak boleh melebihi batas maksimal dua ratus lima puluh kata. Pilihlah kosakata yang representatif untuk digunakan sebagai kata kunci di bagian bawah abstrak, lalu susun kata-kata tersebut secara alfabetis dari huruf A hingga Z dengan jumlah maksimal enam kata kunci yang dipisahkan tanda koma. Pastikan format penulisan kepala abstrak telah sesuai dengan buku pedoman fakultas sebelum dokumen siap untuk dikompilasi ke dalam versi cetak akhir skripsi Anda.
}

% --- Kata Kunci ---
%
%     Maksimal 6 kata, diurutkan dari A sampai Z
%
\newcommand{\kataKunci}{%
	Analisis, \textit{E-Commerce}, Kelayakan, Kualitatif, Skripsi.
}
%------------------------ BATAS AMAN ------------------------

{%
	\setstretch{1}
	
    \begin{center}
        % --- Penulis dan Email Korespondensi ---
        \namaMahasiswa \\
        \texttt{\emailref{\emailKorespondensi}}
        
        % --- Judul ---
        \textbf{\MakeUppercase{\judul}}
        
        % --- Informasi Institusi & Tahun ---
        \daerahKampus, Fakultas \fakultas, Program Studi \prodi, \tahun
    \end{center}
    
    \vspace{0.4cm}
    
    \noindent\abstrak
    
    \vspace{0.4cm}
    
    \noindent\textbf{Kata Kunci:} \kataKunci
}

\newpage